%%%%%%%%%%%%%%%%%%%%%%%%%%%%%%%%%%%%%%%%%
% Imperial College London Poster (Portrait)
% LaTeX Template
% Version 1.0 (February 22, 2024)
%
% For current versions and to report
% issues, please see:
% https://github.com/ImperialCollegeLondon/imperial_latex_templates
%
%!TEX program = xelatex
% Note: this template must be compiled with XeLaTeX rather than PDFLaTeX
% due to the custom fonts used. The line above should ensure this happens
% automatically, but if it doesn't, your LaTeX editor should have a simple toggle
% to switch to using XeLaTeX.
%
% © Imperial College London, 2024. This template, including logo and fonts, is 
% for use of Imperial staff and students only for university business. All rights 
% reserved to the copyright owners.
%
%%%%%%%%%%%%%%%%%%%%%%%%%%%%%%%%%%%%%%%%%

%----------------------------------------------------------------------------------------
%	CLASS, PACKAGES AND OTHER DOCUMENT CONFIGURATIONS
%----------------------------------------------------------------------------------------

\documentclass[
	%print, % Uncomment to convert colors to CMYK for printing purposes
	a1papersize, % Uncomment to use an A1 paper size instead of the default A0 paper size
]{ImperialPoster}

\usepackage[style=numeric,maxnames=2,minnames=2]{biblatex} % bibliography
	\addbibresource{bibliography.bib}
	\AtBeginBibliography{\small}
	\AtEveryBibitem{
		\clearfield{month}
		\clearfield{series}
		\clearfield{venue}
		\clearname{editor}
		\clearlist{publisher}
		\clearlist{location} % alias to field 'address'
		\clearfield{doi}
		\clearfield{url}
		\clearfield{venue}
		\clearfield{issn}
		\clearfield{isbn}
		\clearfield{urldate}
		\clearfield{eventdate}
		\clearfield{pages}
		\clearfield{eventtitle}
		\clearfield{conference}
		\clearfield{journaltitle}
		%\clearfield{number}
		%\clearfield{volume}
	}

\usepackage{subcaption} % subfigures

\usepackage{cleveref}

% https://tex.stackexchange.com/questions/17360/reduce-spacing-in-bibliography-using-biblatex
\usepackage{setspace}

% hacky way to get \pyop2 and \pyop3 as valid macros
% source: https://tex.stackexchange.com/questions/13290/how-to-define-macros-with-numbers-in-them
\def\pyop#1{\ifnum#1=2 {PyOP2}\else \ifnum#1=3 {\texttt{pyop3}}\fi \fi}

\newcommand{\myblue}[1]{\textcolor{ICLBlue}{#1}}

%----------------------------------------------------------------------------------------
%	POSTER INFORMATION
%----------------------------------------------------------------------------------------

\postertitle{\LARGE{\texttt{pyop3:}}\\ \large{An Execution Abstraction for Compact\\ Computational Kernels on Unstructured Meshes}}

% Contributor/author names, can be split across multiple lines manually with \\
% Add affiliations using the \affiliation{} command
\posterauthors{Connor J. Ward\\ David A. Ham\\ Jack D. Betteridge\\~\\c.ward20@imperial.ac.uk}

% Command to output coauthor logos to the right of the Imperial logo in the header of the poster
% If no coauthor logos are required, remove or comment out this command
\coauthorlogos{
	\coauthorlogo[3.5cm]{firedrake.png}
}

%----------------------------------------------------------------------------------------
%	ABSTRACT
%----------------------------------------------------------------------------------------

% The execution of compact computational kernels is ubiquitous in continuum mechanics simulations.
% These "local" kernels, often hand-written, are executed repeatedly as part of an outer loop
% over (usually) the cells or facets of the mesh, with the relevant parts of the global data
% structures loaded into local temporaries as required.
%
% Whilst initially appearing to be a simple task, this pack/unpack execution model is
% surprisingly complicated. Performance portability considerations often require intrusive
% changes (e.g. data layout transformations, different programming languages, manual
% vectorisation) that would require either rewriting the outer loop time and time again or
% sacrificing performance.
%
% In this work we present pyop3, a new framework for handling mesh iteration problems. pyop3
% uses code generation to produce the different loops automatically from a high-level
% representation. Core to its design is a novel data layout abstraction that bridges the gap
% between unstructured mesh data and traditional N-dimensional arrays. Work is currently under
% way to integrate pyop3 into the finite element package Firedrake, with the hope that it
% facilitates the development of new numerical methods and performance optimisations.

%----------------------------------------------------------------------------------------

\begin{document}

\titlesection

\begin{multicols}{2}

%----------------------------------------------------------------------------------------
%	FIRST COLUMN
%----------------------------------------------------------------------------------------

\vspace{-1em}
\section{1. Motivation}

The execution of compact computational kernels, where one loops over a mesh and performs some ``local" calculation for each mesh entity, is ubiquitous in continuum mechanics simulations.
To give an example, in the finite element method one assembles the linear systems to be solved by looping over the cells of the mesh, evaluating an integral calculation using local DoFs, and then adding the result to a global matrix or vector.

As a one-off, this ``outer loop" is very straightforward to write by hand.
However, there is a plethora of design choices (PDE, discretisation, offload device, etc) and possible optimisations (vectorisation, tiling, data layout transformations, etc) that mean that this loop would need to be rewritten time and time again.
This handwritten approach both hampers programmer productivity and inhibits the development of higher level abstractions (e.g.~\cite{alnaesUnifiedFormLanguage2014a}).
\textcolor{ICLBlue}{Therefore, a software abstraction that captures the semantics and automates the construction of ``outer loops" is needed.}

\vspace{-1em}

\section{2. Introducing \pyop3}

To this end, we have written a new library for automating the construction and execution of these loops: \pyop3.
Given some user-provided data structures and a \myblue{loop expression} containing the application of some local kernel, \pyop3 uses \myblue{code generation} to produce low-level code that is near equivalent to something handwritten (\cref{fig:codegen}).

\begin{figure}[H]
	\centering
	\includegraphics[width=.35\textwidth]{codegen.pdf}
	\caption{Simplified example of the code generated from a \pyop3 loop expression.}
	\label{fig:codegen}
\end{figure}

\pyop3 is similar in many respects to its precursor, \pyop2~\cite{rathgeberPyOP2HighLevelFramework2012}, and work to replace \pyop2 with \pyop3 in the finite element framework Firedrake~\cite{FiredrakeUserManual} is ongoing.

\subsection{References}

\vspace{1em}

\begingroup
\setstretch{0.8}
\setlength\bibitemsep{0pt}
\printbibliography[heading=none]
\endgroup

\columnbreak

%----------------------------------------------------------------------------------------
%	SECOND COLUMN
%----------------------------------------------------------------------------------------

\begin{figure}[H]
	\centering
	\hfill
	\begin{subfigure}{.16\textwidth}
		\centering
		\includegraphics[width=.7\textwidth]{ref_triangle.pdf}
	\end{subfigure}
	\begin{subfigure}{.16\textwidth}
		\centering
		\includegraphics[width=.7\textwidth]{lagrange_element_3_vec.pdf}
	\end{subfigure}
	\hfill
	\begin{subfigure}{.16\textwidth}
		\centering
		\includegraphics[width=.7\textwidth]{lagrange_element_2_dg.pdf}
	\end{subfigure}
	\hfill
	% Caption needs to be in a parbox of the same width as the figure in order to span the same width
	\parbox{0.49\textwidth}{
		\caption{
			An example of a mixed finite element (Scott-Vogelius, middle and right) on a reference triangle (left).
			DoFs associated with vertices, edges and cells are shown in green, red and blue respectively.
			The lower-order element (right) is discontinuous, so all the DoFs are considered to belong to the cell.
		}
		\label{fig:sv_element}
	}
\end{figure}

\begin{figure}[H]
	\includegraphics[width=0.49\textwidth]{sv_element_dof_layout.pdf}
	% Caption needs to be in a parbox of the same width as the figure in order to span the same width
	\parbox{0.49\textwidth}{
		\caption{
			A possible data layout for the DoFs of a one-cell mesh for the finite element pair in \cref{fig:sv_element}.
			The labels $V_h$ and $Q_h$ represent the function spaces of the two elements.
		}
		\label{fig:dof_layout}
	}
\end{figure}

\vspace{-1em}

\section{3. A new abstraction for mesh data}

The most significant difference between \pyop2 and \pyop3 is the development of a new abstraction for describing data layouts.
Mesh data has a complicated structure (\cref{fig:dof_layout}) that maps poorly to existing ways of describing data layouts (e.g. N-dimensional arrays).
Whilst \pyop2 handles this structure by throwing away information, \pyop3 is capable of fully describing complex mesh data layouts using \myblue{axis trees} (\cref{fig:axis_tree}).
Having a richer abstraction for data layouts makes composition easier and thus \pyop3 is able to describe and generate code for a wider range of problems than \pyop2.

\begin{figure}[H]
	\includegraphics[width=0.49\textwidth]{axis_tree.pdf}
	% Caption needs to be in a parbox of the same width as the figure in order to span the same width
	\parbox{0.49\textwidth}{
		\caption{
			An axis tree that describes the data layout shown in \cref{fig:dof_layout}.
		}
		\label{fig:axis_tree}
	}
\end{figure}

\vspace{-1em}

\section{4. Future work}

Once \pyop3 has been fully integrated into Firedrake, there is scope for a variety of performance improvements and avenues of research to pursue.
Possible low-level changes include:

\begin{itemize}
	\item GPU support (high priority)
	\item Data layout transformations
	\item Loop tiling
\end{itemize}

Whilst at a high level the following should also become expressible:

\begin{itemize}
	\item Physics-based preconditioners
	\item Slope limiters
	\item Adaptivity
	\item Particle methods
\end{itemize}

%----------------------------------------------------------------------------------------

\end{multicols}

\end{document}
